\documentclass[11pt, reqno]{article}

\usepackage{amsmath, amsthm, amssymb}
\usepackage{enumitem}
\usepackage{tcolorbox}
\usepackage{hyperref}
\usepackage{tikz}
\usepackage{tikz-cd}
\usepackage{pgfplots}
\pgfplotsset{compat=1.18}
\usetikzlibrary{arrows.meta}
\usepackage{mathrsfs}
\usepackage{fancyhdr}
\usepackage[bottom=0.75in, top=1in, left=0.5in, right=0.5in]{geometry}
\usepackage{array}   % for \newcolumntype macro
\newcolumntype{L}{>{$}l<{$}}

\theoremstyle{plain}
\newtheorem*{theorem}{Theorem}
\newtheorem*{proposition}{Proposition}
\newtheorem{exercise}{Exercise}
\newtheorem*{lemma}{Lemma}
\newtheorem*{corollary}{Corollary}

\theoremstyle{definition}
\newtheorem*{definition}{Definition}
\newtheorem*{example}{Example}

\theoremstyle{remark}
\newtheorem*{remark}{Remark}

\renewcommand{\phi}{\varphi}
\renewcommand{\epsilon}{\varepsilon}
\renewcommand{\emptyset}{\varnothing}

\newcommand{\RR}{\mathbb{R}}
\newcommand{\ZZ}{\mathbb{Z}}
\newcommand{\NN}{\mathbb{N}}
\newcommand{\CC}{\mathbb{C}}
\newcommand{\QQ}{\mathbb{Q}}

\DeclareMathOperator{\ima}{\text{im}}

\begin{document}

\topmargin=-40pt
\rhead{Henry Woodburn}
\lhead{Math 656}
\renewcommand{\headrulewidth}{1pt}
\renewcommand{\headsep}{20pt}
\thispagestyle{fancy}

{\Huge \bfseries \noindent Homework 4}
\begin{enumerate}
    \item[1.] Let $T$ be a compact operator. Then $s_j(T)$, the $j$th singular value of $T$,
    is a continuous function of $T$ in the operator norm topology. 

    Suppose $T_n$ is a sequence of compact operators converging to $T$ in the operator norm topology.
    Let $T_n = U_n A_n$ and $T = UA$ be polar decompositions. First, we show that $A_n \to A$ 
    in the operator norm topology as well. We have 
    \[
        \|T_n^* T_n - T^* T\| \leq \|T_n^* - T^*\|\|T_n\| + \|T^*\|\|T_n - T\| \to 0
    \]
    since $\|T^*\| = \|T\|$ and $\|T_n\|$ is bounded since it is a convergent sequence. Then 
    it follows that $A_n = \sqrt{T_n^* T_n} \to \sqrt{T^* T} = A$, since the square root of a positive operator
    is a continuous function in the norm topology. 

    We will apply fischer's principle: Denote the eigenvalues of $A$ in decreasing order 
    by $\alpha_k$, $k = 1, 2, \dots$. Then 
    \[
        \alpha_N = \max_{S_N} \min_{x \in S_N} R_A(x).
    \]

    We have
    \begin{align*}
            \max_{S_N} \min_{x \in S_N} R_A(x) = &\max_{S_N} \min_{x \in S_N} \langle Ax, x\rangle 
            \leq \max_{S_N} \min_{x \in S_N} \langle A_n x, x\rangle + \langle (A - A_n)x, x \rangle \\
            &\leq \|A_n - A\| + \max_{S_N} \min_{x \in S_N} R_{A_n}(x).
    \end{align*}
    By symmetry, we obtain
    \[
        |\max_{S_N} \min_{x \in S_N} R_A(x) - \max_{S_N} \min_{x \in S_N} R_{A_n}(x)| < \|A_n - A\|,
    \]
    and thus the $N$th eigenvalue of $A_n$ converges to that of $A$.

    \item[4.] We prove that the trace class operators form a complete linear space under the trace norm.
    
    Suppose $T_n$ is a cauchy sequence in the trace norm. Then $T_n$ is cauchy in the operator norm,
    since $\|T\| \leq \|T\|_{\text{tr}}$, and thus $T_n$ converges to some bounded operator $T$ 
    in the operator norm. 

    Then we must prove $T$ is trace class, and that $T_n$ converges to $T$ in the trace norm. 

    In our text, it was shown that we can compute the trace norm using the formula 
    \[
        \|T\|_{\text{tr}} = \sum_0^\infty \langle T z_n, w_n\rangle,
    \]
    where $z_n$ is an orthonormal basis for $H$ consisting of eigenvectors of the absolute value $|T|$,
    and $w_n$ is the completion of the orthonormal set $U z_n$ to a basis for all of $H$, where $T = U|T|$.
    Moreover, 
    \[
        \|T\|_{\text{tr}} = \sup \sum_n |\langle T f_n, e_n\rangle|,
    \]

    where the supremum is taken over all pairs of orthonormal bases $\{f_n\}$ and $\{e_n\}$.

    Let $z_n$ and $w_n$ be such bases as above, for the operator norm limit $T$ of the sequence $T_n$. Using the fact
    that operator norm convergence implies strong operator convergence, for any $N > 0$, we can write
    \[
        \sum_1^N |\langle T z_n, w_n\rangle| = \lim_j \sum_1^N |\langle T_j z_n, w_n\rangle|
        \leq \limsup \|T_j\|_{\text{tr}} < \infty,
    \]
    since cauchy sequences are bounded. As this holds for all $N$, $T$ is bounded in the trace norm. 

    To show convergence in the trace norm, let $z_n$ and $w_n$ be as above. Fix $j$. Then 
    \[
        \sum_1^N|\langle (T - T_j) z_n, w_n\rangle|  = \lim_\ell \sum_1^N |\langle(T_\ell - T_j)z_n, w_n\rangle|
        \leq \limsup_\ell \|T_j - T_\ell\|_{\text{tr}}
    \]
    for all $N > 0$. As $\limsup\limits_\ell \|T_j - T_\ell\|$ converges to $0$ in $j$, we have $T_n \to T$ in the 
    trace norm. It follows that the trace class operators are a complete linear space under the trace norm,
    using also that the trace norm satisfies the triangle inequality.

    \item[8.] Let $V$ be the operator on $L^2[0,1]$ defined by integration of a function from $0$ to $s$. 
    The continuity of $Vu(s)$ for $u \in L^2[0,1]$ follows cauchy-schwartz:
    \[
        \left|\int_a^b u(t)dt\right| \leq \|u\|_{L^2}|b - a|.
    \]
    In fact, for $u \in B$, the unit ball of $L^2[0,1]$, we have 
    \[
        \|u\|_{L^2}|b-a| \leq |b-a|,
    \]
    implying equicontinuity of $V(B)$. Together with the fact that every function in $V(B)$ takes 
    the value $0$ at $x = 0$, this implies pointwise boundedness of $V(B)$. 

    By Arzela-Ascoli theorem, $V(B)$ is a compact subset of $C[0,1]$.
    
    \item[9.] Let $V$ be the operator from the previous exercise. Its integral kernel 
    is the function
    \[
        K(s,t) = \begin{cases}
        1 & t < s\\
        0 & t > s
        \end{cases}.
    \]
    Then the integral kernel of the adjoint (in $L^2[0,1]$) is the function
    \[
        K^*(s,t) = K(t,s) = \begin{cases}
            1 & s < t\\
            0 & s > t
        \end{cases}.
    \]

    From this we can deduce that
    \[
        V^* V(u)(s) = \int_s^1 \int_0^r u(t) dt dr.
    \]

    By the fundamental theorem of calculus, we obtain
    \[
        (V^* V(u))'(s) = - \int_0^r u(t)dt
    \]
    and 
    \[
        (V^* V(u))''(s) = -u(t).
    \]
    Thus the inverse of $V^* V$ is the differential operator $-\partial_x^2$. 
    
    By the results of our functional calculus, the eigenvalues of $V^* V$ are squares of those of $|V|$. 
    Denote $A = V^* V$. We seek the values $\lambda$ such that 
    \[
        Au = \lambda u,
    \]
    and thus
    \[
        u = -\lambda \partial_x^2 u.
    \]

    Then we must solve the differential equation
    \[
        u'' + \frac{1}{\lambda}u = 0.
    \]
    We can get the boundary data from
    \[
        0 = (Au)(1) = \lambda u(1),
    \]
    thus $u(1) = 0$, and from
    \[
        0 = (Au)'(0) = \lambda u'(0),
    \]
    thus $u'(0) = 0$. 

    Then we can solve the ODE to obtain solutions $u(t) = \cos(nx)$ for $n = \frac{\pi}{2} + k\pi$ for 
    $k = 0, 1, \dots$, and the eigenvalues of $A$ are of the form
    \[
        \frac{1}{(\frac{\pi}{2} + k\pi)^2}.
    \]

    By the remark above, the singular values are 
    \[
        \frac{1}{\frac{\pi}{2} + k\pi}
    \]
    for $k = 0, 1, \dots$. Since these behave like $\frac{1}{k}$, they are not summable. Thus $V$ is 
    not trace class.

    \item[10.] Suppose $F$ is a pseudoinverse of $G$, with
    \[
        GF = 1 - T,\quad FG = 1 - S,
    \]
    for $S$ and $T$ compact operators. 
    
    Let $\lambda$ be an eigenvalue of $T$, $\lambda \neq 1$, and $e$ be a corresponding eigenvalue. Then
    \[
        Te = \lambda e = e - GFe,
    \]

    so that 
    \begin{equation}\label{one}
        (1 - \lambda)e = GFe.
    \end{equation} Applying $F$ to both sides, we obtain
    \[
        FGFe = (1 - S)Fe = (1- \lambda)Fe,
    \]
    so that $\lambda Fe = SF e$. Then $Fe$ is an eigenvector of $S$ with eigenvalue $\lambda$. This
    shows every eigenvalue of $T$ is one of $S$. Moreover, \ref{one} implies that if $e \neq 0$, 
    then $Fe \neq 0$ as well, and also that if $Fe = 0$, then $e = 0$. Thus the map $x \mapsto Fx$ 
    is injective as a map from the eigenspace of $\lambda$ under $T$ to the eigenspace of $\lambda$ under $S$.

    By symmetry, one has that $G$ is an injective map from the eigenspace of $\lambda$ under $S$
    to the eigenspace of $\lambda$ under $T$.

    From this, we have that the dimension of both eigenspaces are equal and thus the eigenvalues have the 
    same multiplicity. 

\end{enumerate}

\end{document}